\section*{Bab \ref{ch:b3:innate-immunity} --- Imunitas Bawaan dan Peradangan}

\begin{solution}{exo:b3:innate-immunity:1}
\omterm{def:b3:innate-immunity:innate}{Penghadang}: keratin kulit (mekanis), lendir dan \omterm{def:b3:cytoskeleton-motility:cilia}{silium} saluran napas (penjebakan
dan pembersihan), asam lambung (membunuh mikroba yang tertelan), lisozim dan
defensin di dalam sekresi (melisiskan dan melubangi), \omterm{def:b3:microbiomes:symbiosis}{mikrobiom} komensal
(menempati ceruknya). Selnya: \omterm{def:b3:innate-immunity:innate}{neutrofil} (membunuh bakteri lewat \omterm{def:b3:innate-immunity:killing}{fagositosis}),
\omterm{def:b3:innate-immunity:innate}{makrofag} (makan, membersihkan, membunyikan lonceng bahaya), \omterm{def:b3:innate-immunity:innate}{sel dendritik} (mengusung
antigen ke kelenjar limfa lalu memulai tanggapan adaptif), sel mast
(melepaskan histamin, meradangkan), \omterm{def:b3:innate-immunity:innate}{sel pembunuh alami} (membunuh sel terinfeksi dan
sel terubah).
\end{solution}

\begin{solution}{exo:b3:innate-immunity:2}
Dimiliki bersama seluruh kelas mikroba, penting bagi mereka (sehingga ia tidak dapat
dilepaskan untuk lolos), dan tidak ada pada inangnya. \omterm{def:b3:bacteriology:envelope}{Lipopolisakarida} ---
TLR4; flagelin --- TLR5; RNA untai ganda --- TLR3 (dan RIG-I di dalam
sitosol); DNA \omterm{def:b3:chromatin-epigenetics:methylation}{CpG} tak termetilasi --- TLR9; serpihan \omterm{def:b3:bacteriology:envelope}{peptidoglikan} --- protein
NOD; $\beta$-glukan --- dektin-1.
\end{solution}

\begin{solution}{exo:b3:innate-immunity:3}
Kemerahan: pelebaran arteriol oleh histamin, nitrat oksida, dan
prostaglandin. Panas: aliran darah yang naik itu juga. Pembengkakan: kebocoran
plasma pada venula sesudah endoteliumnya mengerut (histamin, C3a, C5a).
Nyeri: bradikinin dan prostaglandin E$_{2}$ yang menajamkan nosiseptor,
ditambah tekanan dari pembengkakannya.
\end{solution}

\begin{solution}{exo:b3:innate-immunity:4}
Klasik: C1q yang mengikat antibodi (atau \omterm{prop:b3:innate-immunity:systemic}{protein C-reaktif}) pada sebuah permukaan.
Lektin: lektin pengikat manosa pada gula mikroba. Alternatif:
denyut sendiri C3 dan pengendapan C3b pada permukaan tak terlindung mana pun.
Hasilnya: C3b mengopsonisasi bagi \omterm{def:b3:innate-immunity:killing}{fagositosis}; C3a dan C5a meradangkan lalu
merekrut; C5b memulai \omterm{def:b3:innate-immunity:complement}{kompleks penyerang membran} yang melisiskan.
\end{solution}

\begin{solution}{exo:b3:innate-immunity:5}
$B = [k_{0}/(a-d)](e^{(a-d)t} - 1)$ dengan $a - d = 0.1$: pada \qty{60}{s},
$20(e^{6} - 1) \approx 8000$; pada \qty{90}{s}, $20(e^{9} - 1) \approx
1.6\times 10^{5}$. Sel inang: $k_{0}/(d - a) = 2/0.1 = 20$.
\end{solution}

\begin{solution}{exo:b3:innate-immunity:6}
Berguling pada \omterm{def:b3:innate-immunity:inflammation}{selektin} (lektin yang mengikat karbohidrat); pengisyaratan
\omterm{def:b3:innate-immunity:inflammation}{kemokin} yang mengaktifkan \omterm{def:b3:innate-immunity:inflammation}{integrin}; pelekatan teguh \omterm{def:b3:innate-immunity:inflammation}{integrin} pada
molekul pelekat endotelium (ICAM); perpindahan menembus antara
sel endotelium; kemotaksis sepanjang landaiannya. Tanpa
\omterm{def:b3:innate-immunity:inflammation}{integrin} $\beta_{2}$, \omterm{def:b3:innate-immunity:innate}{neutrofilnya} berguling tetapi tidak pernah berhenti atau menyeberang:
mereka menumpuk di dalam darah (cacahnya sangat tinggi), infeksinya berulang tanpa
nanah terbentuk, lukanya sembuh buruk, dan tali pusarnya terlepas terlambat.
\end{solution}

\begin{solution}{exo:b3:innate-immunity:7}
Penyakit granulomatosa menahun: \omterm{def:b3:innate-immunity:killing}{oksidase NADPH-nya} cacat, sehingga tidak ada
superoksida, hidrogen peroksida, atau hipoklorit yang dibuat di dalam fagosomnya.
Infeksi oleh jasad yang memusnahkan hidrogen peroksidanya sendiri
(katalase positif): \emph{Staphylococcus aureus}, \emph{Aspergillus},
\emph{Burkholderia}, \emph{Serratia}, \emph{Nocardia}. Granuloma terbentuk
karena \omterm{def:b3:innate-immunity:innate}{makrofag} menelan mikroba yang tidak dapat dibunuhnya lalu menemboknya,
dengan perekrutan terus-menerus, di dalam sebuah lesi menahun.
\end{solution}

\begin{solution}{exo:b3:innate-immunity:8}
Reseptor penghambat mencacah MHC~I; reseptor pengaktif mencacah ligan
tekanan; selnya mati bila pengaktifannya melebihi penghambatannya. (a) Kehilangan
MHC~I menyingkirkan penghambatannya: \omterm{def:b3:innate-immunity:innate}{sel pembunuh alami} membunuhnya --- itulah
ongkos lolos dari sel~T. (b) Dengan MHC~I yang dipertahankan tetapi ligan tekanan
dipajang, pengaktifannya masih dapat menang dan selnya dibunuh; \omterm{def:b3:cancer-biology:hallmarks}{tumor}
karena itu kerap melepaskan atau menekan ligan tekanannya.
\end{solution}

\begin{solution}{exo:b3:innate-immunity:9}
Sebuah protein murni tidak mengusung pola patogen, sehingga \omterm{def:b3:innate-immunity:innate}{sel dendritik} yang
mengambilnya tidak matang, menyajikannya tanpa kostimulasi, dan
sel~T yang mengenalinya dimatikan atau mengabaikannya. Sebuah ajuvan
--- tawas, sebuah lipid, sebuah ligan \omterm{def:b3:innate-immunity:prr}{reseptor serupa Toll} --- memicu reseptor
pola, mematangkan \omterm{def:b3:innate-immunity:innate}{sel dendritiknya}, lalu memasok isyarat yang kedua.
Inilah pokok Janeway: sistem adaptifnya menanggapi antigen hanya
ketika sistem bawaannya mengatakan bahwa antigennya datang bersama sebuah mikroba.
\end{solution}

\begin{solution}{exo:b3:innate-immunity:10}
Tanpa demam, $48$ penggandaan: $10^{4}\times 2^{48} \approx 3\times
10^{18}$ (sebuah kemustahilan yang menunjukkan bahwa pembunuhannya, bukan
pertumbuhannya, yang memutuskan hasilnya). Dengan demam, $32$ penggandaan:
$4\times 10^{13}$ --- yakni faktor $2^{16} \approx 65\,000$ bakteri lebih sedikit
yang dihadapi \omterm{def:b3:innate-immunity:innate}{neutrofilnya}. Ongkosnya: $3\times 0.12\times 8 = \qty{2.9}{MJ}$
sehari, yakni sekali makan besar. Berbaloi bila patogennya diperlambat dan
inangnya punya cadangan; tidak berbaloi bila patogennya tumbuh sama baiknya pada
\qty{40}{\celsius}, atau bila inangnya kelaparan, atau bila demamnya sendiri
mendekati suhu yang berbahaya.
\end{solution}

\begin{solution}{exo:b3:innate-immunity:11}
Mengikat faktor H berarti merekrut pengatur milik inangnya sendiri: $d$ naik di atas
$a$ dan gelungnya meluruh pada mikrobanya seperti pada sebuah sel inang. Memotong
C3b menyingkirkan konvertase yang mengendap: sekali lagi sebuah kenaikan $d$. Sebuah kapsul
menyajikan sebuah permukaan yang di atasnya konvertase terbentuk buruk, yang menurunkan $a$, lalu
menyembunyikan C3b yang mengendap dari reseptor fagosit. Yang termurah:
satu protein permukaan yang mengikat faktor H, yang dipasok inangnya cuma-cuma
dan yang mengerjakan tugasnya.
\end{solution}

\begin{solution}{exo:b3:innate-immunity:12}
Pada sebuah serpihan, pelebaran, kebocoran, dan perekrutannya melibatkan beberapa
mililiter pembuluh lalu menghantarkan pertahanan ke tempat yang memerlukannya; perantara
yang sama yang dilepaskan ke seluruh peredaran melebarkan tiap pembuluh lalu
membocorkan tiap kapiler, sehingga tekanannya runtuh dan pembekuannya
terpicu di mana-mana --- fisiologinya serupa, skalanya
mematikan. TNF mendorong \omterm{def:b3:innate-immunity:inflammation}{peradangan} sinovium pada radang sendi rematoid,
sehingga menghalanginya meredakan penyakitnya; tetapi TNF juga yang menjaga
\omterm{def:b3:innate-immunity:innate}{makrofag} tetap tertata di dalam granuloma yang mengurung basil tuberkel yang
tertidur, dan menghalanginya membiarkan tuberkulosis laten lolos.
\end{solution}

\begin{solution}{pb:b3:innate-immunity:1}
\textbf{1.} Sejam adalah $1.5$ penggandaan: $10^{4}\times 2^{1.5} \approx
2.8\times 10^{4}$.
\textbf{2.} $2\times 10^{4}$ \omterm{def:b3:innate-immunity:innate}{neutrofil} membunuh sampai $4\times 10^{5}$;
bakterinya tumbuh hanya dari $2.8\times 10^{4}$ ke $8\times 10^{4}$:
pembunuhannya melampaui pertumbuhannya lima kali lipat.
\textbf{3.} Jam ke-2: $8\times 10^{4}$ sebelum pembunuhannya, tidak ada sesudahnya.
Populasinya memuncak pada sekitar $8\times 10^{4}$ dekat \qty{2}{h} lalu nol
sejak itu; tidak ada yang tersisa pada \qty{6}{h}.
\textbf{4.} Kedatangan pada \qty{3}{h}: $10^{4}\times 2^{4.5} = 2.3\times
10^{5}$ ketika itu. Jam ke-4: $6.4\times 10^{5} - 4\times 10^{5} = 2.4\times
10^{5}$; jam ke-5: $6.8\times 10^{5} - 4\times 10^{5} = 2.8\times 10^{5}$;
jam ke-6: $7.9\times 10^{5} - 4\times 10^{5} = 3.9\times 10^{5}$ dan
naik --- \omterm{def:b3:innate-immunity:innate}{neutrofilnya} tidak lagi mengimbangi, lalu sebuah bisul terbentuk. Dua
jam penundaan mengubah sebuah kesembuhan menjadi infeksi menahun.
\textbf{5.} $8\times 10^{4}/20 = 4000$ \omterm{def:b3:innate-immunity:innate}{neutrofil} mati sambil membunuh, ditambah
$2\times 10^{4}$ yang tiba lalu mati dengan sendirinya dalam sehari. Nanah adalah
\omterm{def:b3:innate-immunity:innate}{neutrofil} mati, bakteri mati, dan jaringan yang mencair; \omterm{def:b3:innate-immunity:innate}{makrofag} membersihkannya
dengan memakan bangkainya, atau ia harus disalirkan.
\textbf{6.} Kedatangannya $2000$ sejam, membunuh $4\times 10^{4}$: jam ke-2,
$8\times 10^{4} - 4\times 10^{4} = 4\times 10^{4}$; jam ke-3, $1.1\times
10^{5} - 4\times 10^{4} = 7\times 10^{4}$; jam ke-4, $2\times 10^{5} -
4\times 10^{4} = 1.6\times 10^{5}$ --- yang tumbuh tanpa batas. Pasiennya
tidak dapat mengurung sebuah inokulum sepele, sehingga \omterm{def:b3:bacteriology:antibiotics}{antibiotik} harus mengerjakan
pembunuhannya sejak tanda yang pertama.
\textbf{7.} $B = 12.5(e^{0.08t} - 1)$: \qty{30}{s}, $125$; \qty{60}{s},
$1500$; \qty{120}{s}, $1.8\times 10^{5}$.
\textbf{8.} $e^{0.08t} = 1.6\times 10^{5}$: $t = 12/0.08 = \qty{150}{s}$.
\textbf{9.} $k_{0}/(d - a) = 10$.
\textbf{10.} $B(120) = 33(e^{3.6} - 1) \approx 1200$, yakni $150$ kali lebih sedikit.
Kapsulnya membeli waktu semenit --- waktu untuk berlipat ganda, dan perlindungan sampai
antibodi tiba untuk memulai jalur klasik pada kapsulnya
sendiri.
\textbf{11.} $10^{4}\times 2\times 10^{6} = 2\times 10^{10}$ C3b: yakni
sepersejuta C3 plasmanya.
\textbf{12.} Tanpa C3 ketiga jalurnya berhenti pada langkah bersamanya:
tidak ada \omterm{def:b3:innate-immunity:complement}{opsonisasi}, tidak ada \omterm{def:b3:innate-immunity:complement}{anafilatoksin}, tidak ada lisis --- infeksi berulang
yang berat oleh bakteri berkapsul. Tanpa C9 hanya liang terminalnya
yang hilang; \omterm{def:b3:innate-immunity:complement}{opsonisasi} dan \omterm{def:b3:innate-immunity:inflammation}{peradangan} bekerja, dan fenotipenya berupa
infeksi \emph{Neisseria} yang berulang, yakni satu-satunya marga yang buruk
dikendalikan fagosit dan baik dikendalikan lisis.
\textbf{13.} $2\times 0.24\times 8 = \qty{3.8}{MJ}$, yakni sekitar \qty{100}{g}
lemak.
\textbf{14.} Satu penggandaan dalam sejam itu: $2\times 10^{4}$ alih-alih
$2.8\times 10^{4}$, yakni faktor $1.4$.
\textbf{15.} Jam ke-2: $4\times 10^{4}$ sebelum pembunuhannya, tidak ada sesudahnya --- lenyap
pada \qty{2}{h}, seperti tadi; sumbangan demamnya berarti ketika
pasokan \omterm{def:b3:innate-immunity:innate}{neutrofilnya} pas-pasan, seperti pada pertanyaan 4 dan 6.
\textbf{16.} $199\,\text{mg/L}\times 3\,\text{L} = \qty{0.6}{g}$; $0.6/
\num{115000} = 5.2\times 10^{-6}$ mol $= 3.1\times 10^{18}$ molekul;
sepanjang \qty{172800}{s}, yakni $1.8\times 10^{13}$ tiap detik.
\textbf{17.} Obat itu menurunkan kembali titik setelnya ke arah \qty{37}{\celsius},
pasiennya merasa lebih baik, dan bakterinya berganda sedikit lebih cepat; pada
model ini \omterm{def:b3:innate-immunity:innate}{neutrofilnya} tetap menang, dan bukti bahwa obat penurun demam
memperburuk infeksi biasa masihlah lemah.
\textbf{18.} Di atas sekitar \qty{41}{\celsius} protein mulai terdenaturasi,
enzim gagal, dan neuron salah menyala (kejang); \qty{42}{\celsius} adalah
mematikan. Titik setelnya ditahan di bawah itu oleh kriogen --- interleukin-10,
vasopresin, glukokortikoid --- dan oleh langit-langit kerja prostaglandin
di dalam hipotalamus.
\textbf{19.} $10^{9}\times 3\times 10^{6} = 3\times 10^{15}$ molekul
$= 5\times 10^{-9}$ mol di dalam \qty{5}{L}: yakni \qty{1e-9}{mol/L}.
\textbf{20.} Sama dengan kepekatan yang menjenuhkan: tiap \omterm{def:b3:innate-immunity:innate}{makrofag} di
dalam tubuh teraktifkan semaksimalnya sekaligus.
\textbf{21.} $10^{11}\times 1000\times 3600 = 3.6\times 10^{17}$ molekul
$= 6\times 10^{-7}$ mol; di dalam \qty{15}{L}, \qty{4e-8}{mol/L} --- yakni empat
ribu kali kepekatan yang aktif.
\textbf{22.} Tiap arteriol melebar dan tiap venula bocor: volume darahnya
mengumpul di dalam pembuluh yang melebar lalu lolos ke dalam jaringannya, sehingga
tekanannya turun; ginjalnya, yang kurang dialiri, berhenti menapis; dan
pembekuan yang terpicu di dalam ribuan kapiler menyumbatnya lalu memakan habis
faktor pembekuannya.
\textbf{23.} \omterm{def:b3:bacteriology:envelope}{Gram-negatif} yang terbunuh melepaskan seluruh \omterm{def:b3:bacteriology:envelope}{lipopolisakaridanya}
sekaligus, yakni sebuah bolus bagi TLR4, sehingga lonjakan \omterm{def:b3:innate-immunity:inflammation}{sitokinnya} dapat memburuk berjam-jam.
Anti-TNF gagal karena pada saat renjatannya runtunannya telah melewati TNF
--- interleukin-1, interleukin-6, dan pembekuannya sudah berjalan
--- dan karena TNF diperlukan untuk mengendalikan bakterinya: waktunya, bukan
sasarannya, yang menjadi masalah.
\textbf{24.} \omterm{def:b3:bacteriology:envelope}{Lipopolisakarida} yang disuntikkan tidak berbuat apa-apa terhadap
mencit tanpa TLR4, yang selamat dari dosis yang mematikan bagi mencit normal; sebuah infeksi
\omterm{def:b3:bacteriology:envelope}{Gram-negatif} yang hidup membunuhnya, karena ia tidak dapat mendeteksi bakterinya awal lalu
menegakkan \omterm{def:b3:innate-immunity:inflammation}{peradangan} yang mengurungnya. Paradoksnya larut: sang
pendeteksi yang menyebabkan renjatannya adalah pendeteksi yang menyediakan
pertahanannya.
\textbf{25.} Tidak ada bakteri yang tersisa pada \qty{6}{h} (semuanya terbunuh pada \qty{2}{h});
sekitar $1.8\times 10^{5}$ C3b tiap bakteri pada \qty{120}{s}; demam dua hari
berongkos \qty{3.8}{MJ}, yakni sekitar \qty{100}{g} lemak.
\end{solution}
